\section{Go SDK}


\subsection{Package Structure}

\begin{lstlisting}[caption=Go SDK Structure]
hanzo-go/
  hanzo.go
  client.go
  product.go
  order.go
  customer.go
  error.go
  pagination.go
  webhook.go
  go.mod
\end{lstlisting}

\subsection{Client Implementation}

\begin{lstlisting}[language=Go, caption=Go Client]
package main

import (
    "context"
    "fmt"
    "github.com/hanzo.ai/hanzo-go"
)

func main() {
    client := hanzo.NewClient("sk_live_...")

    // List products
    params := &hanzo.ProductListParams{
        Limit: hanzo.Int64(10),
    }
    iter := client.Products.List(params)
    for iter.Next() {
        product := iter.Product()
        fmt.Println(product.Name)
    }
    if err := iter.Err(); err != nil {
        log.Fatal(err)
    }

    // Create a product
    product, err := client.Products.Create(&hanzo.ProductParams{
        Name:        hanzo.String("Widget"),
        Price:       hanzo.Int64(1999),
        Description: hanzo.String("A useful widget"),
    })
    if err != nil {
        log.Fatal(err)
    }
}
\end{lstlisting}

\subsection{Pointer Helpers}

Go requires pointers for optional fields:

\begin{lstlisting}[language=Go, caption=Go Pointer Helpers]
// Helper functions for optional parameters
func String(v string) *string { return &v }
func Int64(v int64) *int64 { return &v }
func Bool(v bool) *bool { return &v }

// Usage
params := &ProductParams{
    Name:  hanzo.String("Widget"),
    Price: hanzo.Int64(1999),
}
\end{lstlisting}

\subsection{Error Handling}

\begin{lstlisting}[language=Go, caption=Go Error Handling]
product, err := client.Products.Create(params)
if err != nil {
    if hanzoErr, ok := err.(*hanzo.Error); ok {
        switch hanzoErr.Code {
        case hanzo.ErrorCodeInvalidRequest:
            fmt.Printf("Invalid param: %s\n", hanzoErr.Param)
        case hanzo.ErrorCodeAuthentication:
            fmt.Println("Invalid API key")
        default:
            fmt.Printf("API error: %s\n", hanzoErr.Message)
        }
    }
    return
}
\end{lstlisting}

